\documentclass{article}

\usepackage{booktabs}
\usepackage{tabularx}
\usepackage{hyperref}

\hypersetup{
    colorlinks=true,       % false: boxed links; true: colored links
    linkcolor=red,          % color of internal links (change box color with linkbordercolor)
    citecolor=green,        % color of links to bibliography
    filecolor=magenta,      % color of file links
    urlcolor=cyan           % color of external links
}

\title{Hazard Analysis\\\progname}

\author{\authname}

\date{}

\input{../Comments.text}
\input{../Common.text}

\begin{document}

\maketitle
\thispagestyle{empty}

~\newpage

\pagenumbering{roman}

\begin{table}[hp]
\caption{Revision History} \label{TblRevisionHistory}
\begin{tabularx}{\textwidth}{llX}
\toprule
\textbf{Date} & \textbf{Developer(s)} & \textbf{Change}\\
\midrule
Date1 & Name(s) & Description of changes\\
Date2 & Name(s) & Description of changes\\
... & ... & ...\\
\bottomrule
\end{tabularx}
\end{table}

~\newpage

\tableofcontents

~\newpage

\pagenumbering{arabic}

\wss{You are free to modify this template.}

\section{Introduction}

\wss{You can include your definition of what a hazard is here.}

\section{Scope and Purpose of Hazard Analysis}

\wss{You should say what \textbf{loss} could be incurred because of the
hazards.}

\section{System Boundaries and Components}

\subsection{Coordinator}
\label{subsec:coordinator}
This component exists on all nodes, but the following responsibilities are only run on the leader node:
\begin{itemize}
	\item Providing an OpenAI-API compatible API to the user for use with tools like
		\href{https://opencode.ai/}{OpenCode}.
	\item Tracking cluster health, and detecting node failures.
	\item Notifying nodes of partitioning changes.
	\item Rebuilding the cluster and saving in-progress requests
		in the case of node failure, removal, or addition.
	\item Measuring node connection performance.
\end{itemize}
On a follower node, it is primarily responsible for managing the inference engine
on its device, following the instructions of the leader node.

\subsection{Model Partitioner}
This component determines how to partition the model among the participating devices. It
decides at which points to cut the graph, how large each partition is, and which device gets
each partition. This component is integral to ensuring that each device is assigned an appropriate
partition given the device's computational and networking performance.

\subsection{Inference Engine}
This component is the high-performance inference engine running on each device. It will handle
execution of the given model partition for each device. It will call into
\href{https://github.com/ggml-org/ggml}{ggml} using the Rust-to-C FFI to control the execution of a
device's model partition.

\subsection{Tensor I/O}
This component handles the compressing, sending, receiving, and decompressing of tensors between devices.
It will leverage \href{https://grpc.io/}{gRPC} for high-performance communication, with the potential
to use a custom solution should \href{https://grpc.io/}{gRPC} fail to meet our requirements in practice.

\subsection{Boundaries}
The system boundaries for Lake primarily consist of:
\begin{itemize}
	\item The \hyperref[subsec:coordinator]{Coordinator's} OpenAI-API compatible API.
	\item Model files (GGUF) loaded from disk.
	\item Networking components between devices (Thunderbolt, Ethernet, Wi-Fi, routers).
	\item The OS and GPU resources on each individual device.
\end{itemize}

\section{Critical Assumptions}

We assume:
\begin{itemize}
	\item The model file (GGUF) provided to Lake is already a functional LLM before partitioning.
	\item User requests are only made through the API on the leader device.
	\item Only one model will be run by the pool at a time.
	\item The local area network connecting the devices is private, and trusted by the user.
\end{itemize}

\section{Failure Mode and Effect Analysis}

\wss{Include your FMEA table here. This is the most important part of this document.}
\wss{The safety requirements in the table do not have to have the prefix SR.
The most important thing is to show traceability to your SRS. You might trace to
requirements you have already written, or you might need to add new
requirements.}
\wss{If no safety requirement can be devised, other mitigation strategies can be
entered in the table, including strategies involving providing additional
documentation, and/or test cases.}

\section{Safety and Security Requirements}

\wss{Newly discovered requirements.  These should also be added to the SRS.  (A
rationale design process how and why to fake it.)}

\section{Roadmap}

\wss{Which safety requirements will be implemented as part of the capstone timeline?
Which requirements will be implemented in the future?}

\newpage{}

\section{Appendix --- Generative AI Use Disclosure}

\wss{ChatGPT (version 5) was used to brainstorm ideas for the template for this disclose section.}

\subsection{AI Tools Used}

\wss{Give the name of the AI tool(s) used and the version.}

\subsection{How and Where Used}

\wss{Explain what the tool(s) were used for.  Examples include: brainstorming,
explaining a technical concept, reviewing the document, generating or modifying
text, generating or modifying diagrams, identifying errors or omissions or
inconsistencies.}

\wss{For each of the uses, identify which parts of the document were affected.}

\wss{You do not need to report every interaction, but you should disclose the 
uses that materially contributed to your work}

\subsection{Verification of AI-Generated Content}

\wss{\begin{itemize}
    \item Describe how the team checked AI-generated or AI-assisted material
    before including it in the document.
    \item In particular, verify factual claims, technical assertions,
    requirements, calculations, references, and other information for which an
    AI system may produce incorrect or fabricated results.
    \item \textbf{The team is responsible for the accuracy, completeness,
consistency, and appropriateness of the submitted document, regardless of
whether material was generated by a human or an AI system.}
\end{itemize}}

\newpage{}

\section{Appendix --- Reflection}

\wss{Not required for CAS 741}

\input{../Reflection.text}

\begin{enumerate}
    \item What went well while writing this deliverable? 
    \item What pain points did you experience during this deliverable, and how
    did you resolve them?
    \item Which of your listed risks had your team thought of before this
    deliverable, and which did you think of while doing this deliverable? For
    the latter ones (ones you thought of while doing the Hazard Analysis), how
    did they come about?
    \item Other than the risk of physical harm (some projects may not have any
    appreciable risks of this form), list at least 2 other types of risk in
    software products. Why are they important to consider?
\end{enumerate}

\end{document}
